\documentclass[11pt]{article}
\usepackage[margin=1in]{geometry}
\usepackage{amsmath,amssymb,amsthm}
\usepackage[hidelinks]{hyperref}

\newtheorem*{theorem}{Literature Answer}

\title{Negative Literature Answer to the RN Continuous-Image Problem in arXiv:1005.4303}
\author{FA/Banach run packet}
\date{2026-06-19}

\begin{document}
\maketitle

\section*{Status}

This is a literature-already-answered packet. It records that the continuous
image problem for Radon--Nikodym compact spaces, stated in the selected-problems
survey of Aviles--Kalenda \cite{AK}, was answered negatively by the later paper
of Aviles--Koszmider \cite{AvKos}. It is not an original proof from this run.

\section*{Original Problem}

In Section 6, ``Radon-Nikodym compact spaces'', Aviles--Kalenda write that the
main open question since Namioka is whether continuous images of RN compacta are
again RN compact. They then ask:

\begin{quote}
Is the continuous image of an RN compact space also RN compact?
\end{quote}

The local source location inspected was
\texttt{data/parsed/arxiv\_sources/1005.4303/probl5.tex}, lines 627--633.

\section*{Supporting Answer}

The supporting paper \cite{AvKos} is titled
\emph{A continuous image of a Radon-Nikodym compact which is not
Radon-Nikodym}. Its abstract states that the authors construct a continuous
image of an RN compact space which is not RN compact, solving Namioka's problem.
The precise theorem used here is Theorem 1.1:

\begin{quote}
There exists a continuous surjection \(\pi:L_0\to L_1\) such that \(L_0\) is
a zero-dimensional RN compact space but \(L_1\) is not RN compact.
\end{quote}

\begin{theorem}
The continuous-image question from Section 6 of \cite{AK} has a negative
answer. A continuous image of an RN compact space need not be RN compact.
\end{theorem}

\begin{proof}[Literature identification]
Theorem 1.1 of \cite{AvKos} supplies a compact space \(L_0\), a compact space
\(L_1\), and a continuous surjection \(\pi:L_0\to L_1\). The domain \(L_0\) is
RN compact, indeed zero-dimensional RN compact, while the image \(L_1\) is not
RN compact. Thus \(L_1\) is a continuous image of an RN compactum but is not RN
compact. This is exactly the negation of the question asked in \cite{AK}.

The supporting authors also explicitly identify the construction as a solution
of the problem posed by Isaac Namioka, so this is an explicit later answer
rather than an agent-identified implication.
\end{proof}

\section*{Scope Limitations}

This packet clears only the headline continuous-image problem at the start of
Section 6 of \cite{AK}. It does not resolve the nearby questions in that survey
about quasi-RN compacta, cardinal reductions, the space of almost increasing
functions, convex hulls of RN compacta, zero-dimensional RN covers, embeddings
into measure spaces over scattered compacta, or unions of two RN compacta.

\section*{Verification Notes}

The source and supporting papers are distinct. The supporting paper appeared in
2011, after the 2010 selected-problems survey, and explicitly says that it solves
Namioka's problem. The local duplicate search covered arXiv ids
\texttt{1005.4303} and \texttt{1112.4152}, the survey title, and phrases
\texttt{continuous image of an RN compact}, \texttt{Radon-Nikodym compact}, and
\texttt{Namioka}.

\begin{thebibliography}{9}

\bibitem{AK}
A. Aviles and O. F. K. Kalenda,
\emph{Compactness in Banach space theory - selected problems},
arXiv:1005.4303.

\bibitem{AvKos}
A. Aviles and P. Koszmider,
\emph{A continuous image of a Radon-Nikodym compact which is not
Radon-Nikodym},
arXiv:1112.4152.

\end{thebibliography}

\end{document}
